%% kapa: exam=1 round=21 year=2010 period=1 form=A subject=economics subject_ko=경제원론 seq=25 number=65
%% source: Q-Net 감정평가사 기출문제 (cst00309) artlSeq=5024181 file=extracted/1008/5024181/2010년 제21회 감정평가사 1차시험 기출문제/1교시.hwp
%% conversion: kordoc 4.24.51 -> tools/gen1.py (auto); review=proofread-sonnet; flags=none
\begin{question}{65}
A국에는 2개의 은행이 있는데, 지급준비율을 제1은행은 20\%, 제2은행은 10\%로 항상 유지한다. 甲은 기존에 보유하고 있던 현금 100만원을 제1은행에 예금하였고, 제1은행은 지급준비금을 제외한 금액을 乙에게 대출하였다. 乙은 이 돈으로 丙에게서 물품을 구입하였고, 丙은 이 대금을 제2은행에 예금하였다. 제2은행은 지급준비금을 제외한 금액을 丁에게 대출하였다. 이상의 거래로부터 추가적으로 창출된 통화량은?
\begin{choices}
  \item 144만원
  \item 152만원
  \item 160만원
  \item 232만원
  \item 332만원
\end{choices}
\end{question}
\answer{2}
